% Appendix A --- included by MAIN_CRYPTOBENCH_GEOAUDIT.tex via \input.
% GENERATED by tools/emit_wire_appendix.py. Do not edit by hand: the
% names, groups, radii and counts are imported from the modules that
% define them, and `make wires` fails if this file drifts from them.

\appendix
\section{Appendix A: every local quantity and every wire}
\label{app:wires}

This appendix is the definition of the input contract. Section~\ref{sec:methods} describes the field in terms of \(43\) local quantities read under a fixed family of neighbourhood statistics; what follows states each of the \(43\), the neighbourhood it is read over and its value at the boundary, and then the rule that expands them into the \NTabWires{} wires. It is generated from the modules, so it cannot describe a version of the code that is not the one that produced the numbers.

\paragraph{Conventions.} \(c_i\) is the centroid of the heavy atoms of residue \(i\) and \(N(i)\) its contacts, meaning residues whose centroid lies within \(10\)~\AA. Neighbourhoods never cross a chain: a residue's neighbourhood is computed inside its own chain, which is what allows one chain to be scored without reference to any other. \(\hat n_i\) is the outward radial unit vector from the chain centroid, \(R_g\) the radius of gyration, \(s_i\) the sequence position, and \(a_i\) the number of heavy atoms. Hydrogens are discarded and a ligand-leak guard rejects any non-solvent heteroatom before any of this is computed.

\paragraph{The probe field.} Quantities in G1 and \texttt{rotor\_gain} are read off a free-space grid: points on a \(1.5\)~\AA\ lattice at least \(2.6\)~\AA\ from every atom, capped at \(6000\) points. The buriedness \(\beta(p)\) of such a point is the fraction of \(30\) Fibonacci-spaced directions along which an atom lies within \(11\)~\AA\ of the ray, at perpendicular tolerance \(1.8\)~\AA.

\paragraph{The ultrametric ball partition.} \(\rho_a\) is the coordination number of atom \(a\) within \(8\)~\AA. Contact edges are deflated to \(\|c_a - c_b\| (\rho_{\max}/\min(\rho_a,\rho_b))^{\gamma}\) with \(\gamma = 2\), and the balls at radius \(\tau\) are the components of the minimum spanning tree restricted to edges of weight at most \(\tau\), which is exactly the closed-ball partition of the minimax ultrametric. \(\tau\) is chosen by the valuation threshold of the tree-weight spectrum, not tuned. Note that the \texttt{valuation} descriptor below uses its own exponent \(0.5\), not \(\gamma\).

\subsection{The \NAlgFeatures{} algebraic and topological quantities}

The grouping is not cosmetic: each group is one dense quaternary table in the cascaded compiler, so no group may exceed six digits.

\paragraph{G1 surface exposure (6).}
\begin{description}
  \item[\texttt{bur}] mean over probe points $p$ attributed to $i$ of the buriedness $\beta(p)$, the fraction of $30$ Fibonacci directions blocked by an atom within $11$~\AA\ of the ray, at perpendicular tolerance $1.8$~\AA. \emph{Read over}: free-grid probe points within $6$~\AA\ of any heavy atom of $i$. \emph{At the boundary}: $0$ when no probe point is attributed to $i$.
  \item[\texttt{void}] $\log\!\left(1 + \#\{p : \beta(p) \ge 0.55\}\right)$. \emph{Read over}: the same attributed probe points. \emph{At the boundary}: $\log 1 = 0$ when none of them is that buried.
  \item[\texttt{probe\_max}] $\max_p \beta(p)$. \emph{Read over}: the same attributed probe points. \emph{At the boundary}: $0$ when no probe point is attributed to $i$.
  \item[\texttt{protrusion}] $\left\| \operatorname{mean}_{j \in N(i)} c_j - c_i \right\| / R$, the offset of the neighbourhood's centroid from the residue. \emph{Read over}: residues within $R = 10$~\AA. \emph{At the boundary}: $0$ when $i$ has no contact.
  \item[\texttt{exposure}] $\#\{p\} / a_i$, probe points per heavy atom of the residue. \emph{Read over}: probe points within $5$~\AA\ of any heavy atom of $i$. \emph{At the boundary}: $0$ when none; $a_i \ge 1$ by construction.
  \item[\texttt{concavity}] $\operatorname{mean}_{j \in N(i)} \widehat{(c_j - c_i)} \cdot \hat n_i$ with $\hat n_i$ the outward radial unit vector from the chain centroid. \emph{Read over}: residues within $10$~\AA. \emph{At the boundary}: $0$ when $i$ has no contact.
\end{description}

\paragraph{G2 local spectral geometry (6).}
\begin{description}
  \item[\texttt{fiedler}] $\lambda_2$ of the combinatorial Laplacian $L = D - A$ of the lining $G[S_i]$, $S_i = \{i\} \cup N(i)$. \emph{Read over}: induced subgraph on $i$ and its contacts, capped at the $48$ nearest. \emph{At the boundary}: $0$ when $i$ has no contact or $|S_i| \le 1$.
  \item[\texttt{lap\_max}] $\lambda_{\max}$ of the same $L$. \emph{Read over}: the same lining. \emph{At the boundary}: $0$ when $i$ has no contact.
  \item[\texttt{spec\_gap}] $\lambda_2 / \max(\lambda_{\max}, 10^{-9})$. \emph{Read over}: the same lining. \emph{At the boundary}: the floor in the denominator is what makes it defined at $\lambda_{\max} = 0$.
  \item[\texttt{mean\_degree}] $2E / |S_i|$ of the lining. \emph{Read over}: the same lining. \emph{At the boundary}: $0$ when $i$ has no contact.
  \item[\texttt{betti0}] number of connected components of $G[S_i \setminus \{i\}]$ --- how many disjoint walls face the residue. \emph{Read over}: the lining re-thresholded at $7$~\AA, because at $10$~\AA\ the lining of a globular protein is connected and this count is the constant $1$. \emph{At the boundary}: $0$ when the lining is empty.
  \item[\texttt{betti1}] cycle rank $E - V + C$ of the same graph, an integer. \emph{Read over}: the same $7$~\AA\ lining. \emph{At the boundary}: $0$ when the lining is empty.
\end{description}

\paragraph{G3 density field calculus (6).}
\begin{description}
  \item[\texttt{rho}] mean over the residue's heavy atoms of the atomic packing density $\rho$, the coordination number within $8$~\AA. \emph{Read over}: the residue's own atoms. \emph{At the boundary}: $\rho \ge 1$ by construction.
  \item[\texttt{grad\_rho}] $\|g_i\|$ where $g_i$ solves the normal equations $A_i g_i = b_i$ of the first-order Taylor fit, $A_i = \sum_j (c_j - c_i)(c_j - c_i)^\top$, $b_i = \sum_j (\rho_j - \rho_i)(c_j - c_i)$. \emph{Read over}: residues within $10$~\AA. \emph{At the boundary}: $g_i = 0$ when $A_i$ is singular, which is a neighbourhood spanning fewer than three dimensions.
  \item[\texttt{grad\_radial}] $g_i \cdot \hat n_i$, the outward component of the same gradient. \emph{Read over}: the same neighbourhood. \emph{At the boundary}: $0$ with $g_i$.
  \item[\texttt{lap\_rho}] $\operatorname{mean}_{j \in N(i)} \rho_j - \rho_i$, the graph Laplacian of the density field. \emph{Read over}: residues within $10$~\AA. \emph{At the boundary}: $0$ when $i$ has no contact.
  \item[\texttt{var\_rho}] $\max\!\left(\operatorname{mean}_j \rho_j^2 - (\operatorname{mean}_j \rho_j)^2,\, 0\right)$. \emph{Read over}: the same neighbourhood. \emph{At the boundary}: clipped at zero so rounding cannot produce a negative variance.
  \item[\texttt{rho\_rank\_ball}] rank of $\rho_i$ among the residues of $i$'s ultrametric ball, scaled to $[0,1]$. \emph{Read over}: the ball, not a radius. \emph{At the boundary}: $0.5$ when the ball holds one residue, which is the value that asserts nothing.
\end{description}

\paragraph{G4 ultrametric structure (6).}
\begin{description}
  \item[\texttt{loose}] mean over the residue's atoms of $\bar\rho_{B(a)} / \bar\rho$, the atom's ball density against the global mean. \emph{Read over}: the residue's atoms and their balls. \emph{At the boundary}: $\bar\rho$ is floored at $1$.
  \item[\texttt{ball\_size}] $\log(1 + |B(i)|)$ in atoms. \emph{Read over}: the ball. \emph{At the boundary}: $\log 1 = 0$ for an empty ball, which cannot occur.
  \item[\texttt{ball\_radius}] $\tau \cdot \bar\rho / \bar\rho_{B(i)}$, the ball's cut radius deflated by its own density. \emph{Read over}: the ball. \emph{At the boundary}: the denominator is floored at $10^{-9}$.
  \item[\texttt{ball\_interface}] fraction of $i$'s contacts lying in a different ball. \emph{Read over}: residues within $10$~\AA. \emph{At the boundary}: $0$ when $i$ has no contact.
  \item[\texttt{valuation}] $\max_{j \in N(i)} \|c_j - c_i\| \left(\rho_{\max} / \min(\rho_i, \rho_j)\right)^{0.5} \big/ \tau$: the longest density-deflated edge the residue carries, in units of the global ball radius. \emph{Read over}: residues within $10$~\AA. \emph{At the boundary}: $0$ when $i$ has no contact; the deflator's denominator is floored.
  \item[\texttt{rotor\_gain}] largest increase in probe buriedness produced by rotating one ultrametric ball rigidly by $\theta = 0.20$~rad about its branch axis. \emph{Read over}: the $8$ largest balls of at least $8$ atoms that are not the whole structure; probes within $19$~\AA\ of the hinge centre. \emph{At the boundary}: $0$ when no ball qualifies, which is the honest reading that the structure has no hinge at this scale.
\end{description}

\paragraph{G5 curvature and anisotropy (6).}
\begin{description}
  \item[\texttt{anisotropy}] $(\lambda_1 - \lambda_3)/\operatorname{tr} M$ for $M_i = \operatorname{mean}_{j \in N(i)} (c_j - c_i)(c_j - c_i)^\top$ with $\lambda_1 \ge \lambda_2 \ge \lambda_3$. \emph{Read over}: residues within $10$~\AA. \emph{At the boundary}: $\operatorname{tr} M$ floored at $10^{-12}$.
  \item[\texttt{planarity}] $(\lambda_2 - \lambda_3)/\lambda_1$. \emph{Read over}: the same tensor. \emph{At the boundary}: $\lambda_1$ floored at $10^{-12}$.
  \item[\texttt{sphericity}] $\lambda_3/\lambda_1$. \emph{Read over}: the same tensor. \emph{At the boundary}: $\lambda_1$ floored at $10^{-12}$.
  \item[\texttt{normal\_div}] $\operatorname{mean}_{j} (\hat n_j - \hat n_i) \cdot (c_j - c_i)$, the discrete divergence of the outward normal field. \emph{Read over}: residues within $10$~\AA. \emph{At the boundary}: $0$ when $i$ has no contact.
  \item[\texttt{angle\_deficit}] $1 - \operatorname{mean}_j \angle(\hat n_i, \hat n_j) / \pi$, bounded in $[0,1]$ by construction. \emph{Read over}: residues within $10$~\AA. \emph{At the boundary}: $1$ when $i$ has no contact, the value a flat isolated patch takes.
  \item[\texttt{hess\_bur}] $\operatorname{mean}_{j} \mathrm{bur}_j - \mathrm{bur}_i$, the graph Laplacian of the buriedness field. \emph{Read over}: residues within $10$~\AA. \emph{At the boundary}: $0$ when $i$ has no contact.
\end{description}

\paragraph{G6 global position (5).}
\begin{description}
  \item[\texttt{depth}] $\|c_i - \bar c\| / R_g$ against the chain centroid and radius of gyration. \emph{Read over}: the whole chain. \emph{At the boundary}: $R_g$ floored at $1$.
  \item[\texttt{coordination}] $|N(i)|$, an integer. \emph{Read over}: residues within $10$~\AA. \emph{At the boundary}: $0$ is attainable and meaningful.
  \item[\texttt{contact\_order}] $\operatorname{mean}_{j \in N(i)} |s_j - s_i| / n$, the mean sequence separation of the contacts normalised by chain length. \emph{Read over}: residues within $10$~\AA. \emph{At the boundary}: $0$ when $i$ has no contact.
  \item[\texttt{ball\_offset}] $\|c_i - \bar c_{B(i)}\| / R_g$, displacement from the residue's own ball centroid. \emph{Read over}: the ball. \emph{At the boundary}: $R_g$ floored at $1$.
  \item[\texttt{depth\_rank}] rank of $\mathrm{depth}$ within the chain, scaled to $[0,1]$. \emph{Read over}: the whole chain. \emph{At the boundary}: $0$ for a chain of one residue.
\end{description}

\subsection{The seven physicochemical constants and the propensity}

None of these is fitted to CryptoBench. The seven are published constants of the residue type; the eighth is a count.

\begin{description}
  \item[\texttt{kd}] Kyte--Doolittle hydropathy. \emph{Source}: Kyte and Doolittle, \emph{J.\ Mol.\ Biol.}\ \textbf{157}:105 (1982).
  \item[\texttt{volume}] side-chain volume in \AA$^3$. \emph{Source}: Zamyatnin, \emph{Prog.\ Biophys.\ Mol.\ Biol.}\ \textbf{24}:107 (1974).
  \item[\texttt{aromatic}] aromatic rings in the side chain: 1 for F, Y, H, 2 for W, 0 otherwise. \emph{Source}: chemical structure.
  \item[\texttt{charge}] formal charge at pH 7: $-1$ for D, E, $+1$ for K, R, $+0.1$ for H as partially protonated, 0 otherwise. \emph{Source}: chemical structure.
  \item[\texttt{hbd}] side-chain hydrogen-bond donors. \emph{Source}: chemical structure.
  \item[\texttt{hba}] side-chain hydrogen-bond acceptors. \emph{Source}: chemical structure.
  \item[\texttt{chi}] rotatable side-chain bonds ($\chi$ angles). \emph{Source}: chemical structure.
\end{description}

An unknown or non-standard residue takes the mean of the twenty standard values rather than a sentinel, so it is neutral in the rank order rather than extreme in it.

\paragraph{The propensity, and why it is training-only.} The forty-third quantity is \(\hat P(\text{cryptic} \mid \text{residue type})\), Laplace-smoothed by one pseudo-count per class,
\begin{equation}
  \pi_t = \frac{1 + \sum_{i : \mathrm{type}(i) = t} y_i}{2 + \#\{i : \mathrm{type}(i) = t\}},
\end{equation}
counted over the training partition only. It is a bincount, not a fit. The table is written into the compiled artifact and read back at inference, so a test residue never contributes a count to the table that scores it; a type absent from training receives the neutral prior \(1/2\) rather than an undefined ratio. This is the only one of the \(43\) that touches a label, and it is the reason the field is described as compiled on the training partition rather than as training-free.

\subsection{From \(43\) local quantities to \NTabWires{} wires}

Each local quantity \(x\) is read under five kinds of statistic over intra-chain neighbourhoods \(N_r(i)\) of fixed radius:
\begin{align}
  \text{identity} &: x_i, \\
  \text{mean} &: \textstyle \frac{1}{|N_r(i)|} \sum_{j \in N_r(i)} x_j, \\
  \text{dispersion} &: \Bigl( \textstyle \frac{1}{|N_r(i)|} \sum_j x_j^2 - \bigl(\frac{1}{|N_r(i)|}\sum_j x_j\bigr)^2 \Bigr)^{1/2}, \\
  \text{centred difference} &: x_i - \textstyle \frac{1}{|N_r(i)|}\sum_j x_j, \\
  \text{local rank} &: \textstyle \frac{1}{|N_r(i)|} \bigl|\{ j \in N_r(i) : x_j < x_i \}\bigr|.
\end{align}

The radii are not the same for every statistic, and the counts are what produce \NTabWires{}:

\begin{center}
\begin{tabular}{llr}
\toprule
statistic & radii (\AA) & groups \\
\midrule
identity & --- & 1 \\
mean & 6, 10, 14, 20, 26 & 5 \\
dispersion & 6, 14, 20 & 3 \\
centred difference & 6, 14, 20 & 3 \\
local rank & 6, 14, 20 & 3 \\
\midrule
total & & 15 \\
\bottomrule
\end{tabular}
\end{center}

So \(35\) algebraic quantities, \(7\) constants and one propensity give \(43\) local quantities, and \(43 \times 15 = 645\) wires. The local rank is the only statistic in the set that is not a moment: it asks where a residue sits in the order of its neighbourhood rather than how far it is from the neighbourhood's centre, so it is unchanged by any monotone rescaling of the quantity, and it separates a residue that is marginally the most hydrophobic of its neighbours from one that is marginally the least --- a distinction the centred difference blurs into two small numbers of opposite sign.

\paragraph{Banding.} Each wire is then banded into a quaternary digit by its rank \emph{within its own chain}, ties sharing a mid-rank:
\begin{equation}
  d_i = \min\!\left( \left\lfloor \frac{r_i}{\max(n-1,1)} \, L \right\rfloor,\, L - 1 \right), \qquad L = 4 .
\end{equation}
Ranking inside the chain is what makes a wire mean the same thing in a 57-residue chain and a 307-residue one, and it removes any dependence on absolute units, so no constant has to be carried from the training fold in order to score a new structure.

